\subsection{対称式を使う}\label{節：基本対称式型}
先ほどの\bookterm{recurrence}{漸化式}を考える.
\[
  a_1=3, \quad a_{n+1}=a_n^2-2
\]
邪魔な$-2$を消すには,二乗の展開から$2$を生み出したい.そこで,
\[
  (a+b)^2=a^2+2ab+b^2
\]
の$2ab$に注目する.$ab=1$なら,ここからちょうど$2$が生まれる.
    そこで積が$1$になる二数として,$a=\alpha,b=\dfrac{1}{\alpha}\ (\alpha\ne0)$と選ぶ.
\[
  \left(\alpha+\frac{1}{\alpha}\right)^2=\alpha^2+2+\frac{1}{\alpha^2}
\]
から,\textbf{二乗して2を引くと,二つのべきの指数が2倍になる}.
第$n$項の指数は$2^{n-1}$.あとは初項に合う$\alpha$を選ぶ.

\noindent \bookblocklink{example-symmetric-1}{problem}{answer}{【問題】}\par\nobreak
\begin{smallquote}
  次の\bookterm{recurrence}{漸化式}で定義される\bookterm{sequence}{数列}の\bookterm{general}{一般項}を求めよ.
  \[
    a_1=3, \quad a_{n+1}=a_n^2-2
  \]
\end{smallquote}
\bookblocklink{example-symmetric-1}{answer}{problem}{【解答】}
  \begin{smallquote}
  \bookterm{general}{一般項}が
    \[
      b_n=\alpha^{2^{n-1}}+\alpha^{-2^{n-1}}
    \]
    で表される\bookterm{sequence}{数列}$b_n$を考える.
    \begin{align*}
      b_{n+1}&=\alpha^{2^{n}}+\alpha^{-2^{n}}\\
      &=\left(\alpha^{2^{n-1}}+\alpha^{-2^{n-1}}\right)^2-2\\
      &=b_n^2-2
    \end{align*}
    と更新が一致する.初項を合わせると,
    \begin{align*}
      &\alpha+\alpha^{-1}=3\\
      &\alpha^2-3\alpha+1=0\\
      &\therefore \quad \alpha=\frac{3\pm \sqrt{5}}{2}
    \end{align*}
    二つの根は互いに逆数で,同じ$b_n$を与える.初項と更新の一致から,
    \[
      a_n=\left(\frac{3+ \sqrt{5}}{2}\right)^{2^{n-1}}+\left(\frac{3- \sqrt{5}}{2}\right)^{2^{n-1}}
    \]
  \end{smallquote}

ここで,定数を $-2$ から $-1$ に変えてみよう.
\[
  a_{n+1}=a_n^2-1
\]
同じ置換では,
\begin{align*}
  a_n^2-1
    &=\left(u+u^{-1}\right)^2-1\\
    &=u^2+u^{-2}+1,
\end{align*}
と余分な1が残り,同じ形で二乗を繰り返すことができない.
係数が近いだけでは同じ方法を使えず,恒等式との一致が必要なのである.

同じ考え方は,三次式にも現れる.
\[
  \left(\alpha+\frac{1}{\alpha}\right)^3
  -3\left(\alpha+\frac{1}{\alpha}\right)
  =\alpha^3+\frac{1}{\alpha^3}
\]
対称な組み合わせを保ったまま,指数を3倍できる.

和を取るか差を取るかで,打ち消せる項の符号も変わる.
差の三乗を展開すれば,
\[
  \left(\alpha-\frac{1}{\alpha}\right)^3
  +3\left(\alpha-\frac{1}{\alpha}\right)
  =\alpha^3-\frac{1}{\alpha^3}
\]
となる.指数$3^{n-1}$の候補を例題で確かめよう.

\noindent\bookblocklink{example-symmetric-2}{problem}{answer}{【問題】}\par\nobreak
\begin{smallquote}
  $\displaystyle a_1=\frac{8}{3},\quad a_{n+1}=a_n^3+3a_n$
\end{smallquote}
\bookblocklink{example-symmetric-2}{answer}{problem}{【解答】}
\begin{smallquote}
  \bookterm{general}{一般項}が
    \[
      b_n=\alpha^{3^{n-1}}-\alpha^{-3^{n-1}}
    \]
    で表される\bookterm{sequence}{数列}$b_n$を考える.
    \begin{align*}
      b_{n+1}&=\alpha^{3^{n}}-\alpha^{-3^{n}}\\
      &=b_n^3+3b_n
    \end{align*}
    と更新が一致する.初項を合わせると,
    \begin{align*}
      &\alpha-\alpha^{-1}=\frac{8}{3}\\
      &3\alpha^2-8\alpha-3=0\\
      &\therefore \quad \alpha=3,\,-\frac13
    \end{align*}
    $3^{n-1}$は奇数なので,どちらを選んでも$b_n$は同じ値になる.
    初項と更新の一致から,$\alpha=3$を選べば,
    \[
      a_n=3^{3^{n-1}}-3^{-3^{n-1}}
    \]
\end{smallquote}
